%!TEX root = hems_proj.tex
%%%%%%%%%%%%%%%%%%%%%%%%%%%%%%%%%%%%%%%%%%%%%%%%%%%%%%%%%%%%%%%%%%%%%%%
\chapter{A projekt módosított változatának jelenlegi állapota}\label{ch:FRISSITES}
%%%%%%%%%%%%%%%%%%%%%%%%%%%%%%%%%%%%%%%%%%%%%%%%%%%%%%%%%%%%%%%%%%%%%%%

\begin{osszefoglal}
Ez a fejezet az eredeti dolgozatot kiegészítve összefoglalja a módosított projektverzió fő újdonságait: a valós árintegrációt, az új output-struktúrát, a bővített vizualizációkat és a jelenlegi benchmark-beállításokat.
\end{osszefoglal}

\section{Mi változott az eredeti projekthez képest}
A módosított verzió az alap kutatási kérdést nem változtatja meg: továbbra is természetből inspirált algoritmusokat hasonlítunk össze HEMS optimalizálási feladaton. A különbség a bemenetek és a riportolás realitásában jelentkezik:
\begin{itemize}
  \item valós inverteres PV idősor használata napokra bontva,
  \item valós ENTSO-E day-ahead árak beépítése (RO zóna, órás felbontás),
  \item rendezett kimeneti mappastruktúra,
  \item részletesebb napi és összegző PNG ábrák.
\end{itemize}

\section{Új bemenetek és adatfolyam}
A benchmark futásban kiválasztott napok száma 16. A valós PV profilok az inverter-adatokból származnak, az árprofilok pedig ENTSO-E CSV forrásból kerülnek feldolgozásra. A futtatás során a rendszer minden naphoz 24 elemű órás vektort használ (load, PV, price). A jelenlegi konfigurációban:
\begin{itemize}
  \item \texttt{LOAD\_REPLICAS = 1},
  \item \texttt{RUNS\_PER\_PROFILE = 2},
  \item 5 algoritmus kerül kiértékelésre (GA(P=50), GA(P=100), PSO, GWO, Hybrid).
\end{itemize}

Az adatfolyam implementációs sorrendje:
\begin{enumerate}
  \item inverteres idősorokból órás PV előállítás (napi bontás),
  \item ár-CSV feldolgozás és RON/kWh konverzió,
  \item napi profilok összeépítése (
  \texttt{real\_pv} és \texttt{synthetic}),
  \item algoritmusfuttatás és nyers eredmények mentése,
  \item aggregálás, gap/winner számítás, report és ábragenerálás.
\end{enumerate}

\section{Árbetöltés és fallback logika}
A valós árbetöltésnél először a \texttt{GUI\_ENERGY\_PRICES\_*.csv} állományok közül a legnagyobb (tipikusan éves) fájl kerül kiválasztásra. A feldolgozás során területi szűrés (RO zóna), órás újramintavételezés és hiányérték-kezelés történik. Ha egy napra nincs teljes 24 órás valós árvektor, a rendszer TOU fallback árat használ.

Ez a megközelítés két szempontból fontos:
\begin{itemize}
  \item robusztus futtatást biztosít hiányos adat esetén is,
  \item transzparensen jelzi, mely napoknál történt fallback.
\end{itemize}

\section{Összegző metrikák és győzteskép}
A frissített kiértékelés nem csak átlagköltséget számol, hanem profile-szintű versenyképességi mutatókat is. A rendszer minden profilnál meghatározza az algoritmusonkénti legjobb költséget, majd ehhez viszonyítja a globális napi legjobbat. Ebből képződik az átlagos résmutató (gap), valamint a győzelmi arány.

Ez a riportolási forma jobban elkülöníti:
\begin{itemize}
  \item a stabilan jó algoritmust (alacsony átlagos gap),
  \item az alkalmi nyertes algoritmust (magas, de szórt győzelmek),
  \item és a gyors, de gyengébb minőségű megoldásokat (runtime előny, költség-hátrány).
\end{itemize}

\section{Reprodukálhatóság és kimeneti artefaktok}
A pipeline minden lépése determinisztikus seedeléssel futtatható, így ugyanazon inputok mellett a nyers és aggregált kimenetek reprodukálhatók. Az artefaktok szétválasztása (\texttt{tables}, \texttt{reports}, \texttt{plots/days}, \texttt{plots/summary}) támogatja a tudományos visszakövethetőséget:
\begin{itemize}
  \item \texttt{tables}: adatforrás a statisztikai ellenőrzéshez,
  \item \texttt{reports}: emberileg olvasható összegzés,
  \item \texttt{plots}: kvalitatív és kvantitatív vizuális validáció.
\end{itemize}

Ennek köszönhetően a dolgozatban szereplő következtetések közvetlenül visszavezethetők a mentett nyers mérési sorokra.

\section{Kimeneti struktúra}
A korábbi ``ömlesztett'' mentés helyett a projekt kimenetei rendezett formában készülnek:
\begin{itemize}
  \item \texttt{ekim2339\_proj/outputs/tables}: nyers és aggregált CSV/JSON eredmények,
  \item \texttt{ekim2339\_proj/outputs/reports}: markdown riport,
  \item \texttt{ekim2339\_proj/outputs/plots/days}: napi PV/load/net/ár ábrák,
  \item \texttt{ekim2339\_proj/outputs/plots/summary}: összegző ábrák (cost, gap, runtime, scatter, napi trend).
\end{itemize}

\section{Példa ábrák a módosított futásból}
\begin{figure}[H]
\centering
\includegraphics[width=0.92\textwidth]{outputs/plots/days/2025-06-12_realpv_price.png}
\caption{Napi részletes ábra: valós PV, generált fogyasztás, nettó terhelés és valós órás árprofil}
\label{fig:updated_day_plot}
\end{figure}

\begin{figure}[H]
\centering
\includegraphics[width=0.82\textwidth]{outputs/plots/summary/real_pv_mean_cost.png}
\caption{Összegző ábra: átlagos költség algoritmusonként (real\_pv dataset)}
\label{fig:updated_summary_cost}
\end{figure}

\begin{figure}[H]
\centering
\includegraphics[width=0.92\textwidth]{outputs/plots/summary/real_pv_cost_scatter_runs.png}
\caption{Run-szintű pontfelhő: költségeloszlás algoritmusonként (real\_pv)}
\label{fig:updated_scatter}
\end{figure}

\section{Aktuális eredménykép és értelmezés}
A módosított pipeline jelenlegi futása szerint mindkét dataseten a GA(P=100) adta a legstabilabb összképet átlagos eltérés (gap) alapján. Ugyanakkor napfüggően továbbra is előfordul, hogy más algoritmus nyer egyes profilokon. Ez megerősíti, hogy a probléma továbbra is szcenárió-érzékeny, tehát egyetlen algoritmus abszolút dominanciája nem feltételezhető minden körülmény között.

\section{Következő fejlesztési lépések}
A projekt következő mérnöki és kutatási lépései:
\begin{enumerate}
  \item külön vételi/eladási ármodell bevezetése,
  \item akkumulátor degradációs költség expliciten beépítése a célfüggvénybe,
  \item opcionálisan több fogyasztási profil-változat (replika) ugyanarra a napra a robusztusság növeléséhez,
  \item érzékenységvizsgálat különböző degradációs és árparaméterekre.
\end{enumerate}

A fenti bővítések az eredeti dolgozat irányát megtartják, de a gyakorlati alkalmazhatóságot és a riportolhatóságot jelentősen erősítik.